\section{Military custody and the responsibility gap}
\label{sec:16.4}
In Gaza, the approval stage supplied authorization without judgment. In Maven, the vendor supplies the infrastructure making unprecedented throughput possible while locating responsibility entirely within the military organization. At machine tempo the reviewer can be treated as a component, the supplier can stand outside the decision, and the institution can remain responsible only in the abstract.

Military custody is where the chapter's problem is most acute: the operator is the sovereign that can compel the institutions meant to constrain it — classify the evidence, define the reviewer's jurisdiction, terminate a contract, restrict a supplier, and decide whether an adverse finding reaches anyone able to stop the operation.

A targeting decision passes through more parties than the chain of command names: the commander who authorizes, the operator who uses, the contractor who integrates (whose embedded personnel may sit close enough to the targeting cell that “effectively acting as a commander” is a real question), the commercial vendor whose model supplies part of the capability and whose personnel are nowhere near it, and the system producing the output. The questions that matter are which of them could have known how the system was used, which had authority to refuse or stop it, and which can be made to answer afterward.

The American domestic rule most concerned with human judgment, DoD Directive 3000.09, requires autonomous and semi-autonomous weapon systems to allow “appropriate levels of human judgment over the use of force” \autocite{DoD2023Directive300009}. A system that generates and ranks targets without firing anything is decision support, which plausibly falls outside it. Maven sits in that gap: the domestic rule written for human judgment doesn't cover the system that compressed it.

Article 28 of the Rome Statute \autocite{icc1998rome} is the doctrine built for the case where the person who authorized an act didn't perform it — a commander is responsible for crimes by forces under effective control where they knew \emph{or should have known} and failed to take necessary measures. Neither branch asks what a subordinate intended. But the Court's jurisdiction runs to natural persons, and the provision presupposes forces that can commit a crime and be disciplined — a targeting system is an instrument. So the question is whether the human elements can still be proved once the instrument is doing the work.

\runin{The fourth party no doctrine covers}

The operator, the command and the weapon's developer are all inside the military. The model's developer isn't. Command responsibility doesn't extend to the vendor — not because it sits outside a chain of command (Article 28's civilian branch was written for superiors with no chain of command) but because nobody at the company has effective authority over the military personnel who commit the underlying crime, and an airstrike isn't an activity within any vendor employee's effective responsibility. A use policy, a licensing term, a right to refuse service, even a technical cut-off are not that relationship. Individual liability under Article 25 stays available in principle with demanding mental elements — facilitation requires acting \emph{for the purpose} of facilitating — and ordinary supply of a general-purpose model approaches neither.

Product liability is the other half and it's untested, with two obstacles: the model isn't defective in the sense the doctrine most readily covers, and a sophisticated intermediary chose the application. Where the doctrine has been rewritten for software it moved away from this case, as the EU's directive shows below. Untested isn't exempt, but no vendor is currently answering to it. The vendor sits in a gap nobody designed, produced by applying to it two bodies of law written for soldiers and for manufacturers.

There is also a jurisdictional gap alongside the doctrinal one: ICC jurisdiction is territorial or by nationality, so one of the two campaigns here sits inside a live instrument and the other sits outside it.

\runin{“I do not know whether my system was used.”}

Matthias's responsibility gap is control distributed until nobody holds enough of it to answer \autocite{matthias2004responsibility} — a gap in \emph{causal control}. What the model developer's case exhibits is different and more tractable: \emph{a gap in causal visibility.} Asked about his company's models and Minab, Anthropic's chief executive said, “we don't know exactly how these models were used” \autocite{tnw2026amodeicircuit}. Take the answer as true. The question is how an arrangement was built in which it could be true.

The paperwork for four July 2025 frontier-model agreements, released in August 2026 under a Freedom of Information Act suit \autocite{lasst2026dodai,biddle2026contracts}, shows something more exact than an absence of oversight terms. One amended agreement carries a section headed System Oversight, withheld in full; the released documents for two of the other three carry no such section at all.

And whether a developer reserved an inspection right matters little. A right exercised inside an accredited classified environment produces findings classified with it. Public accountability needs a body that can initiate scrutiny without belonging to the state and say afterwards what it found. No commercial term supplies that, and the redaction is the demonstration: whatever the withheld section says, the people the deployment is aimed at can't read it.

Matthias's gap may not be closable, but absent visibility is an instrumentation choice: logging is solved, contractual inspection rights are ordinary, and a vendor that wanted to know could require, as a condition of licensing into a targeting platform, that categories of use be recorded and made available to a named auditor.

The government's own acquisition guidance is the better evidence, because it's complete. OMB's April 2025 memorandum lists what agencies should secure from vendors, from data rights and portability to pricing transparency, and every item runs from the vendor to the government; none runs back \autocite{omb2025aiacquisition}. The contracting vehicle explains the absence. The four agreements were \emph{other transactions}, which carry no FAR provisions by default, so a contracting officer decides case by case which protections to write in, and GAO has found these instruments risk reduced accountability and transparency \autocite{gao2026aiacquisitions}. A System Oversight section can be missing without anyone having removed it.

The pressure that matters is the advantage attached to removing judgment from the timeline. If one state fields systems that act faster because nobody can intervene before engagement, competitors acquire a reason to match the speed and give up the check. “Meaningful human control” stays true as a description of formal authority while becoming false as a description of what anyone can do before the system acts. And the pressure reaches the reviewers: a military treating delay as operational risk has reason to narrow what requires clearance and move review from authorization-before to explanation-after. Nothing has formally abolished oversight; it has been moved to after the act. So the enforceable question is who can slow or suspend the system without first getting permission from the command whose advantage depends on keeping it fast.
